\documentclass[10pt,conference]{IEEEtran}
\usepackage{cite}
\usepackage{amsmath,amssymb,amsfonts}
\usepackage{algorithmic}
\usepackage{graphicx}
\usepackage{textcomp}
\usepackage{xcolor}
\usepackage{booktabs}
\usepackage{url}
\usepackage{hyperref}

\def\BibTeX{{\rm B\kern-.05em{\sc i\kern-.025em b}\kern-.08em
    T\kern-.1667em\lower.7ex\hbox{E}\kern-.125emX}}

\begin{document}

\title{recovery\_eval: State-Matched and Provenance-Aware Evaluation of Recovery Behavior in Language-Model Reasoning}

\author{\IEEEauthorblockN{Sham Satish Thakare}
\IEEEauthorblockA{\textit{Independent Researcher} \\
Pune, Maharashtra, India \\
Email: shamthakare3000@gmail.com}
}

\maketitle

\begin{abstract}
Evaluating whether post-training procedures enhance a language model's ability to recover from intermediate reasoning errors remains challenging due to confounding variables in state trajectory comparisons. Naive benchmark comparisons often conflate overall accuracy gains with recovery-specific capability. In this work, we introduce \texttt{recovery\_eval}, a reproducible data-centric evaluation framework for machine learning models that pairs verifier-defined error states with prospectively matched reference control states using frozen structural covariates. Our framework incorporates strict evidence governance, including append-only exposure ledgers, primitive neural-rollout provenance, and independently verifiable reconstruction. We demonstrate the framework across 400 genuine continuations generated by two released checkpoints (\texttt{Qwen2.5-Math-1.5B} Base and Instruct) on 20 fresh GSM8K evaluation items (mean normalized weighted-L1 distance $d_{\text{mean}} = 0.0360$, $d_{\max} = 0.0360$). Under the state-matched protocol, we observed continuation success improvements for Instruct over Base in both recovery ($+0.4300$) and control ($+0.5400$) states, yielding a matched recovery-specific contrast of $D_{\text{recovery}} = -0.1100$ with a 95\% descriptive problem-level bootstrap interval of $[-0.240, +0.030]$. These findings demonstrate that aggregate checkpoint gains do not automatically translate into a detectable recovery-specific advantage, highlighting the utility of state-matched evaluation infrastructure for reasoning model diagnostics.
\end{abstract}

\begin{IEEEkeywords}
Language Models, Mathematical Reasoning, Error Recovery, State-Matched Evaluation, Data-Centric AI, Benchmark Governance, Reproducibility.
\end{IEEEkeywords}

\section{Introduction}
Large language models (LLMs) trained on mathematical reasoning tasks exhibit improved step-by-step problem-solving capabilities \cite{cobbe2021gsm8k, qwen25math2024}. A central research question in machine learning evaluation is whether post-training mechanisms (e.g., instruction tuning, process supervision, or reinforcement learning) instill structural intelligence that enables models to detect and recover from early arithmetic or logical missteps during trajectory generation \cite{lightman2023process, zelikman2022star}.

However, evaluating error recovery behavior is susceptible to statistical confounding. Simply comparing model performance on error-containing prefixes against valid prefixes conflates trajectory depth, remaining solution length, problem difficulty, and token complexity. Aggregate accuracy improvements on benchmarks can mask whether a post-trained checkpoint genuinely possesses superior recovery mechanisms or merely exhibits higher baseline fluency across all states.

To address these challenges, we present \texttt{recovery\_eval}, a data-centric evaluation framework designed to isolate recovery-specific behavior in language-model reasoning trajectories. The primary contributions of this paper are threefold:
\begin{enumerate}
    \item \textbf{State-Matched Evaluation Protocol}: A verifier-defined recovery and control state evaluation protocol that enforces prospective structural matching on intermediate reasoning prefixes.
    \item \textbf{Provenance \& Exposure-Governance Architecture}: An infrastructure that preserves primitive BPE token-level neural rollout evidence, weight manifests, and append-only exposure ledgers for independent reconstruction.
    \item \textbf{Framework Demonstration}: A genuine two-checkpoint \texttt{Qwen2.5-Math-1.5B} demonstration showing that aggregate checkpoint accuracy gains ($+0.4300$ recovery vs $+0.5400$ control) do not automatically correspond to a detectable recovery-specific advantage ($D_{\text{recovery}} = -0.1100$).
\end{enumerate}

\section{Related Work}
Recent work has investigated self-correction, test-time compute scaling, and backtracking in reasoning LLMs \cite{snell2024scaling, wang2022selfconsistency}. Table \ref{tab:literature} compares \texttt{recovery\_eval} with existing evaluation paradigms.

\begin{table}[htbp]
\caption{Comparison with Existing Reasoning Evaluation Frameworks}
\label{tab:literature}
\centering
\begin{tabular}{p{2.2cm}p{1.4cm}p{1.6cm}p{1.8cm}}
\toprule
\textbf{Framework / Work} & \textbf{State Matching} & \textbf{Exposure Ledger} & \textbf{Primitive Provenance} \\
\midrule
End-to-End GSM8K \cite{cobbe2021gsm8k} & None & No & Final Answer Only \\
PRM800K \cite{lightman2023process} & Unmatched & Partial & Step Scores \\
STaR \cite{zelikman2022star} & Unmatched & No & Filtered Traces \\
\textbf{\texttt{recovery\_eval} (Ours)} & \textbf{Prospectively Matched} & \textbf{Append-Only LEDGER} & \textbf{Full BPE Token + Raw JSONL} \\
\bottomrule
\end{tabular}
\end{table}

Unlike prior benchmarks that evaluate full trajectories or un-matched intermediate states, \texttt{recovery\_eval} enforces prospective matching on intermediate reasoning prefixes prior to rollout generation.

\section{Problem Definition}
Let $s$ denote a partial reasoning trajectory prefix for a problem instance $p \in \mathcal{P}$. A deterministic verifier $\mathcal{V}(s)$ evaluates prefix validity. 
A \textit{recovery state} $s_R$ contains a verifier-identified intermediate error. A \textit{reference control state} $s_C$ represents a valid intermediate prefix for the same problem at an equivalent problem-solving stage.

For a target policy $\pi$ and baseline policy $\pi_0$, let $V(s) \in \{0, 1\}$ denote the binary continuation success indicator of completing the solution correctly from prefix $s$. We define the matched recovery-specific contrast $D_{\text{recovery}}$ as:
\begin{equation}
D_{\text{recovery}} = \mathbb{E}[V_{\pi}(s_R) - V_{\pi_0}(s_R)] - \mathbb{E}[V_{\pi}(s_C) - V_{\pi_0}(s_C)]
\end{equation}

\section{The \texttt{recovery\_eval} Framework}
The \texttt{recovery\_eval} package provides an end-to-end modular pipeline for state construction, matching, execution, and provenance logging.

\subsection{Recovery and Control State Construction}
Recovery states $s_R$ are constructed by introducing controlled single-step arithmetic perturbations into valid intermediate reasoning prefixes. Reference control states $s_C$ maintain valid prefix steps.

\subsection{Prospective Matching Protocol}
To ensure statistical comparability, each recovery state $s_R$ is matched to a control state $s_C$ using a normalized weighted L1 Manhattan distance over 6 pre-group structural covariates:
\begin{equation}
d(i, j) = \sum_{k=1}^K w_k \frac{|x_{ik} - x_{jk}|}{s_k}
\end{equation}
subject to exact categorical matching on reasoning operation type and problem difficulty. Table \ref{tab:covariates} lists the prospective matching specification.

\begin{table}[htbp]
\caption{Frozen Prospective Matching Covariates and Scales}
\label{tab:covariates}
\centering
\begin{tabular}{lccc}
\toprule
\textbf{Covariate Name} & \textbf{Type} & \textbf{Weight ($w_k$)} & \textbf{Scale ($s_k$)} \\
\midrule
Trajectory Depth & Continuous & 0.4 & 1.5 \\
Remaining Solution Length & Continuous & 0.4 & 1.0 \\
Token Length & Continuous & 0.2 & 15.0 \\
Reasoning Operation & Categorical & Exact & -- \\
Problem Difficulty & Categorical & Exact & -- \\
\bottomrule
\end{tabular}
\end{table}

\subsection{Exposure \& Provenance Governance}
All evaluation items are managed by an append-only event ledger using SHA-256 parent hash chaining to prevent data leakage or post-hoc item selection. Raw rollouts store full input token IDs, generated BPE token IDs, monotonic generation timestamps, model weight manifest hashes, and verifier outputs.

\section{Experimental Validation}
We validate the framework on 20 fresh, prospectively isolated GSM8K test items ($N=20$). For each item, 1 matched recovery state and 1 matched control state are evaluated across 2 model configurations (\texttt{Qwen2.5-Math-1.5B} Base and Instruct) using 5 stochastic generation seeds ($S=5$), yielding exactly 400 rollouts ($20 \times 2 \times 2 \times 5 = 400$).

\begin{table}[htbp]
\caption{Evaluation Design and Provenance Statistics}
\label{tab:provenance}
\centering
\begin{tabular}{lp{5.2cm}}
\toprule
\textbf{Attribute} & \textbf{Specification / Value} \\
\midrule
Base Model & \texttt{Qwen/Qwen2.5-Math-1.5B} (\texttt{4a83ca6e}) \\
Instruct Model & \texttt{Qwen/Qwen2.5-Math-1.5B-Instruct} (\texttt{aafeb0fc}) \\
Hardware Device & Apple Silicon MPS (\texttt{mps:0}) \\
Execution Mode & Non-interactive PyTorch \texttt{model.generate()} \\
Total Rollouts & 400 (200 Base, 200 Instruct) \\
Total Generated Tokens & 19,212 BPE tokens \\
Measured Duration & 1,755.86 seconds ($\approx 29.26$ minutes) \\
Base Throughput & 11.86 tokens/sec (11,354 tokens in 957.37s) \\
Instruct Throughput & 9.84 tokens/sec (7,858 tokens in 798.49s) \\
BPE Decode Match & 100.0\% (400/400 exact match) \\
Raw JSONL SHA-256 & \texttt{51b5a157d9e44102caeb86d0b356f558...} \\
\bottomrule
\end{tabular}
\end{table}

\section{Empirical Results}
Table \ref{tab:outcomes} summarizes continuation success rates and the resulting contrast $D_{\text{recovery}}$.

\begin{table}[htbp]
\caption{Continuation Outcomes and Matched Recovery Contrast}
\label{tab:outcomes}
\centering
\begin{tabular}{lccc}
\toprule
\textbf{State Condition} & \textbf{Base ($\pi_0$)} & \textbf{Instruct ($\pi$)} & \textbf{Difference ($\Delta$)} \\
\midrule
Recovery States ($s_R$) & 0.1500 & 0.5800 & $+0.4300$ \\
Control States ($s_C$) & 0.3800 & 0.9200 & $+0.5400$ \\
\midrule
\textbf{Matched Contrast ($D_{\text{recovery}}$)} & -- & -- & \textbf{$-0.1100$} \\
\bottomrule
\end{tabular}
\end{table}

Under the evaluated state-matched protocol, we did not observe evidence of a recovery-specific advantage for the Instruct checkpoint over the Base checkpoint. The estimated matched recovery-specific checkpoint-interface contrast was $-0.1100$, with a 95\% descriptive problem-level bootstrap interval of $[-0.240, +0.030]$ ($10,000$ resamples).

Notably, the Instruct checkpoint exhibited higher continuation success than Base across both recovery ($+0.4300$) and control ($+0.5400$) conditions. However, because the gain was larger for control states, the net contrast $D_{\text{recovery}}$ is negative. This demonstrates that aggregate post-training accuracy gains should not be interpreted automatically as recovery-specific improvement.

\section{Ablations and Matching Sensitivity}
We evaluate matching quality across the 20 matched pairs. Mean normalized weighted-L1 distance is $d_{\text{mean}} = 0.0360$, median is $d_{\text{median}} = 0.0360$, with maximum distance $d_{\max} = 0.0360$. All 20 pairs (20/20) satisfy both the standard matching threshold ($d \le 0.25$) and tight threshold ($d \le 0.10$). Per-covariate Standardized Mean Differences (SMDs) are reported in Table \ref{tab:sensitivity}.

\begin{table}[htbp]
\caption{Matching Distance, Covariate SMDs \& Threshold Sensitivity}
\label{tab:sensitivity}
\centering
\begin{tabular}{lccc}
\toprule
\textbf{Metric / Covariate} & \textbf{Raw Difference} & \textbf{SMD / Distance} & \textbf{Status} \\
\midrule
Trajectory Depth & $+0.0000$ & $+0.0000$ & Exact Match \\
Remaining Length & $+0.0000$ & $+0.0000$ & Exact Match \\
Token Length & $+2.0000$ & $+0.1333$ & Balanced \\
\midrule
Mean Pair Distance ($d_{\text{mean}}$) & -- & $0.0360$ & $20/20 \le 0.25$ \\
Max Pair Distance ($d_{\max}$) & -- & $0.0360$ & $20/20 \le 0.10$ \\
\bottomrule
\end{tabular}
\end{table}

\section{Limitations}
This evaluation has explicit scope boundaries:
\begin{itemize}
    \item \textbf{Model Scope}: Evaluated on a single model family (\texttt{Qwen2.5-Math-1.5B}) across two released checkpoints.
    \item \textbf{Benchmark Scope}: Evaluation uses 20 GSM8K test items ($N=20$). Benchmark pretraining contamination cannot be ruled out.
    \item \textbf{Sample Size}: 5 stochastic continuations per state/policy serve as evaluation samples, not independent training replications.
    \item \textbf{No Causal Claim}: Results represent a descriptive contrast under controlled state matching, not a causal claim regarding post-training mechanisms.
\end{itemize}

\section{Reproducibility \& Artifact Availability}
All code, raw evidence, and verification scripts are publicly available. The sealed raw rollout dataset \texttt{RAW\_NEURAL\_ROLLOUTS.jsonl} (SHA-256: \texttt{51b5a157d9e44102caeb86d0b356f558aa7499f6bad3634f668f0dd1ed76b1b4}) is archived under Git commit \texttt{2252cf13adb4b929d4b85ffc909e8ea9089ba041}. The publication certificate \texttt{PUBLICATION\_EMPIRICAL\_CERTIFICATE\_V2.json} (SHA-256: \texttt{3f3291ab...}) is committed under \texttt{8228f1c0}.

\section{Conclusion}
We introduced \texttt{recovery\_eval}, a state-matched evaluation framework for language-model reasoning trajectories. By prospectively matching recovery states with control states and maintaining strict evidence governance, the framework enables fine-grained diagnostics of post-training behavior.

\bibliographystyle{IEEEtran}
\bibliography{references}

\end{document}
